\scanpage{4}
\uline{Corollaire} 2 \quad L'hom $K(B/A)\longrightarrow K(B_0/A_0)$ définit un isomorphisme
\[
H^i(B/A)\xrightarrow{\sim}H^i(B_0/A_0)\qquad\text{pour }i\ne0
\]
et une suite exacte
\[
0\longrightarrow I\longrightarrow H^0(B/A)\longrightarrow H^0(B_0/A_0)\longrightarrow0
\]
% margin note written sideways
\margin{d'où $\operatorname{Coker}(A\longrightarrow H^0(B/A))\simeq\operatorname{Coker}(A_0\longrightarrow H^0(B_0/A_0))$}

\inserted{\uline{Cor. 3} a) Pour que $H^i(B/A)=0$ il f. et s. que $H^i(B_0/A_0)=0$ ($i\ne0$ donné)

\deleted{\unsure{b.}} b) $A\longrightarrow\deleted{\unsure{B}}H^0(B/A)$ inj (surj) sss $A_0\longrightarrow H^0(B_0/A_0)$ inj (surj)}

Plus généralement, tensorisant par un $A$-module $M$ plat, on trouve

\uline{Corollaire} \deleted{\unsure{3}}4 \deleted{\unsure{bis}} \quad Soit $M$ un $A$-module \uuline{plat}. On a une suite exacte de \deleted{complexes} \inserted{$A$-modules simpliciaux}
\[
0\longrightarrow C_{M\otimes_A I}\longrightarrow K(B/A,M)\longrightarrow K(B_0/A_0,M_0)\longrightarrow0
\]
d'où des isomorphismes
\[
H^i(B/A,M)\simeq H^i(B_0/A_0,M_0)\qquad i\ne0
\]
et une suite exacte
\[
0\longrightarrow \unsure{\underline M|M}\otimes_A I\longrightarrow H^0(B/A,M)\longrightarrow H^0(B_0/A_0,M_0)\longrightarrow0
\]
% archive sheet number: 3
